%========================
% chapters/02_foundations.tex
%========================
\chapter{Foundations}
\label{chap:foundations}

This chapter introduces the concepts needed to understand the dissertation. The presentation is intentionally centered on the connection between graph topology and sequential generation, because the main object of study is not graph generation in the abstract, but the way in which a graph becomes a sequence before it is given to an autoregressive model.

\section{Graphs, Isomorphism, and Node Orderings}
\label{sec:graphs-orderings}

A graph is denoted by $G=(V,E)$, where $V$ is a finite set of nodes and $E\subseteq V\times V$ is a set of edges. In this dissertation, the experimental graphs are undirected and unweighted, although the sequential generation framework can be extended to richer settings. The number of nodes is $n=|V|$, and the number of edges is $m=|E|$. An adjacency matrix $A\in\{0,1\}^{n\times n}$ represents the same graph only after an ordering of the nodes has been chosen. If the nodes are permuted, the adjacency matrix changes even though the abstract graph remains the same.

This distinction is central. Two graphs $G=(V,E)$ and $G'=(V',E')$ are isomorphic when there exists a bijection $f:V\rightarrow V'$ such that $(u,v)\in E$ if and only if $(f(u),f(v))\in E'$. A generative model that learns graphs should ideally assign the same structural meaning to all isomorphic labelings. However, a neural model trained on matrices or sequences typically observes a particular labeling. Node ordering is therefore the bridge between the invariant object and the ordered representation used for learning.

For a graph $G$, let $\mathcal{O}_G$ be the set of all admissible node orderings. An ordering $\sigma=(v_1,\ldots,v_n)$ is a permutation of $V$. Once $\sigma$ is fixed, the graph can be translated into a sequence of construction actions. In DGMG, for example, the position of a node in $\sigma$ determines when it is introduced, which previous nodes can be selected as edge destinations, and when the model must decide to stop adding edges for the current node. Thus, ordering does not only relabel the graph; it changes the temporal structure of the supervised learning problem.

Ordering affects the generator through three mechanisms. The first is locality. If the correct edge destination is usually close in the ordering, then the generator can rely on recent construction context. The second is entropy. If an ordering makes one destination much more plausible than the others, the action distribution becomes sharper and easier to learn. The third is regularity. If similar graphs induce similar action patterns, training examples reinforce one another; if equivalent graphs induce highly variable sequences, the model must learn many incompatible histories.

These mechanisms explain why the dissertation does not treat graph ordering as a mere preprocessing operation. A fixed ordering rule already encodes a hypothesis about which partial graphs should be easy to continue. A learned ordering policy extends this idea by making the hypothesis trainable and graph-specific.

The dissertation studies several families of graphs because each one emphasizes a different structural constraint. Cycles require every node to have degree two and require a single ring-like closure. Trees require connectivity and acyclicity. Binary trees add a branching constraint. Stars are dominated by a single hub connected to all leaves. Barab\'asi--Albert graphs are sparse scale-free graphs generated by growth and preferential attachment \cite{barabasi1999emergence}. These families were chosen because their topology makes the effect of ordering interpretable.

\section{Autoregressive Graph Generation}
\label{sec:autoregressive-generation}

An autoregressive graph generator decomposes graph generation into a sequence of conditional decisions. Rather than sampling a whole graph at once, it samples actions $a_1,\ldots,a_L$ and defines
\begin{equation}
    p_{\phi}(a_1,\ldots,a_L)
    = \prod_{t=1}^{L} p_{\phi}(a_t \mid a_{<t}),
\end{equation}
where $\phi$ denotes the generator parameters and $a_{<t}$ is the history of previous actions. In graph generation, the history corresponds to a partial graph. This decomposition is attractive because it mirrors a constructive process: nodes and edges can be added incrementally, and validity constraints can be checked during generation.

The price of this decomposition is sensitivity to the chosen action sequence. If a graph can be serialized in multiple ways, each serialization produces a different factorization of the same underlying object. For a graph $G$ and ordering $\sigma$, let $\mathcal{T}(G,\sigma)$ denote the translator that produces a sequence $a=(a_1,\ldots,a_L)$. The negative log-likelihood optimized during teacher-forced training is
\begin{equation}
    \mathcal{L}_{\mathrm{NLL}}(\phi;G,\sigma)
    = -\sum_{t=1}^{L}\log p_{\phi}(a_t \mid a_{<t}),
    \qquad a=\mathcal{T}(G,\sigma).
\end{equation}
Changing $\sigma$ changes $a$ and therefore changes the conditional prediction task. In this sense, ordering is a structural inductive bias imposed before the generator is trained.

\section{Graph Neural Networks}
\label{sec:gnn-foundations}

Graph Neural Networks encode graph structure by repeatedly exchanging messages along edges \cite{scarselli2009gnn,gilmer2017mpnn}. A general message-passing layer updates a node representation by aggregating information from its neighborhood:
\begin{equation}
    h_v^{(k+1)} = \psi^{(k)}\left(h_v^{(k)},\; \square_{u\in\mathcal{N}(v)}\, \eta^{(k)}(h_v^{(k)},h_u^{(k)},e_{uv})\right),
\end{equation}
where $h_v^{(k)}$ is the representation of node $v$ at layer $k$, $\mathcal{N}(v)$ is the neighborhood of $v$, $\eta$ is a message function, $\square$ is a permutation-invariant aggregation operator such as sum or mean, and $\psi$ is an update function. Graph Convolutional Networks simplify this update through normalized neighborhood aggregation \cite{kipf2017gcn}, while Graph Attention Networks use learned attention weights to let the model assign different importance to different neighbors \cite{velickovic2018gat}.

GNNs are important in this dissertation for two reasons. First, autoregressive models uses graph representations to condition the next construction decision. Second, the learned-ordering policy must observe the entire input graph and choose a node to place next in the ordering. A policy that ignores the graph would reduce to a generic permutation sampler, whereas a GNN-based policy can condition its choices on degree, neighborhood structure, frontier effects, and global graph signals.

\section{Reinforcement Learning and Policy Gradients}
\label{sec:rl-foundations}

Reinforcement learning studies sequential decision problems in which an agent interacts with an environment and receives rewards \cite{sutton2018reinforcement}. In a Markov Decision Process, the agent observes a state $s_t$, samples an action $a_t$ from a policy $\pi_\theta(a_t\mid s_t)$, moves to a new state, and eventually receives a return. The objective is to maximize expected return:
\begin{equation}
    J(\theta)=\mathbb{E}_{\tau\sim\pi_\theta}[R(\tau)],
\end{equation}
where $\tau$ is a trajectory and $R(\tau)$ is its return.

When actions are discrete, the reward is delayed, and the environment is not differentiable, policy-gradient methods are a natural choice. The REINFORCE estimator \cite{williams1992reinforce} uses the log-derivative identity to write
\begin{equation}
    \nabla_\theta J(\theta)
    =\mathbb{E}_{\tau\sim\pi_\theta}\left[R(\tau)\nabla_\theta\log p_\theta(\tau)\right].
\end{equation}
In the node-ordering problem, the trajectory is the sequence of selected nodes, and the reward is computed only after the resulting ordering has been used to train and evaluate the graph generator. This makes the problem a direct instance of delayed-reward structured decision-making.

%\section{Why Ordering Changes the Learning Problem}
%\label{sec:ordering-learning-problem}

%Ordering affects the generator through three mechanisms. The first is locality. If the correct edge destination is usually close in the ordering, then the generator can rely on recent construction context. The second is entropy. If an ordering makes one destination much more plausible than the others, the action distribution becomes sharper and easier to learn. The third is regularity. If similar graphs induce similar action patterns, training examples reinforce one another; if equivalent graphs induce highly variable sequences, the model must learn many incompatible histories.

%These mechanisms explain why the dissertation does not treat graph ordering as a mere preprocessing operation. A fixed ordering rule already encodes a hypothesis about which partial graphs should be easy to continue. A learned ordering policy extends this idea by making the hypothesis trainable and graph-specific.
